\begingroup
\renewcommand{\arraystretch}{1.45}
\Needspace{45\baselineskip}
\InterpolationArticle{27}
Thus, for the case where the number of proposed values of the function $X$ forming a complete period is a composite number $=\pi=\mu\nu$, we have learned in Articles 25, 26 to derive, by dividing that period into $\nu$ periods of $\mu$ terms, the same function satisfying all the given values as would result from applying the general theory directly to the entire period. Anyone who tries it in practice will learn, however, that the former method reduces the tedium of mechanical calculation more. No further explanation will now be needed of how this division can be extended still further and applied when the total number of proposed values is a number composed of three or more factors, for example, if the number $\mu$
\SourceLeaf{317}
were again composite, in which case each period of $\mu$ terms can evidently be subdivided into several smaller periods. Of the many other remarks useful for practical application, we shall touch here only on the following. The values of sums such as
\[
\begin{aligned}
&A\cos\lambda a+B\cos\lambda b+C\cos\lambda c\\[5pt]
&\qquad{}+D\cos\lambda d+\text{etc.}=p\\[5pt]
&A\sin\lambda a+B\sin\lambda b+C\sin\lambda c\\[5pt]
&\qquad{}+D\sin\lambda d+\text{etc.}=q
\end{aligned}
\]
(where, as before, we suppose $b-a=c-b=d\quad c\ \text{etc.}=\tfrac1\mu360^{\mathord{0}}=\Delta$)
can often, if $a=0$ does not yet hold, be determined more conveniently by the formulas
\[
\begin{aligned}
&A+B\cos\lambda\Delta+C\cos2\lambda\Delta\\[5pt]
&\qquad{}+D\cos3\lambda\Delta+\text{etc.}=P\\[5pt]
&B\sin\lambda\Delta+C\sin2\lambda\Delta\\[5pt]
&\qquad{}+D\sin3\lambda\Delta+\text{etc.}=Q
\end{aligned}
\]
$p=P\cos a-Q\sin a$, $q=P\sin a+Q\cos a$, or alternatively, on setting
$\frac QP=\operatorname{tang}\varphi$ and $\frac P{\cos\varphi}=\frac Q{\sin\varphi}=R$,
\hfill\break{}by the formulas $p=R\cos(\varphi+a)$, $q=R\sin(\varphi+a)$.

\endgroup
